\documentclass[12pt]{article}
\usepackage[utf8]{inputenc}
\usepackage[T1]{fontenc}
\usepackage{hyperref}
\title{Lolopo Language Resources for K-12 Curriculum Development}
\author{Compiled by Claude Sonnet 5.5 (claude-sonnet-5-5) for Andr\'as Kornai}
\date{}
\begin{document}
\maketitle

\section{Introduction}
Lolopo (ISO 639-3 ycl; Glottocode lolo1259) is a Loloish (Sino-Tibetan) language spoken by roughly half a million Yi people, mainly in central Yunnan, China. Its two main dialects are Lolopo (Nanhua) and Lipo (Dayao), both in Chuxiong Yi Autonomous Prefecture. The web search found general reference pages and one related dictionary site. It found no Lolopo textbook, primer or school programme. This report is therefore short.

\section{Primary Educational Resources}
No Lolopo primers, textbooks or school materials turned up in the search. Lolopo is not known to be taught in schools.

\section{Linguistic and Reference Materials}
\begin{itemize}
\item Wikipedia, Lolopo language (overview of dialects and classification): \url{https://en.wikipedia.org/wiki/Loloish_languages} and \url{https://en.wikipedia.org/wiki/Lolopo_language}
\item Lipo numerals page, Max Planck Institute for Evolutionary Anthropology: \url{https://lingweb.eva.mpg.de/channumerals/Lipo.htm}
\item Webonary ``Mayou'' dictionary. It covers the Yao'an dialect, which the site calls Loxrlavu. This may be a different variety from the Lolopo of ycl, so check before relying on it: \url{https://www.webonary.work/mayou/overview/introduction/}
\end{itemize}
No grammar of Lolopo itself was found.

\section{Religious and Community Resources}
None were found in the search.

\section{Digital and Audio Resources}
\begin{itemize}
\item Omniglot page on Lolopo (Loxrlapor), with writing-system information: \url{https://omniglot.com/writing/lolopo.htm}
\end{itemize}
No audio resources were found.

\section{Curriculum Development Context}
Very little material exists for building a Lolopo curriculum. Any programme would need new primers and orthography decisions, working with the Yi community and with the existing Yi writing systems in China. These resources are reference background only.
\end{document}